\section{Experiments}

We organize the evaluation around eight questions. Section~IV-A compares physical estimators and
general-purpose language decoders, then examines measured inference cost. Section~IV-B tests
generalization across frequency samplings, antenna arrays, spatial blocks, and unseen maps.
Section~IV-C studies robustness to complex AWGN, Section~IV-D isolates semantic alignment and physical
prediction-head components, and Section~IV-E analyzes the dominant LoS--NLoS failure modes.
Section~IV-F evaluates uncertainty-based selective prediction, while Section~IV-G validates the delay
pipeline on measured SISO data. Section~IV-H evaluates whether the learned representation and predicted
physical attributes transfer to downstream localization, cyclic-prefix selection, outage prediction,
MCS adaptation, and transmit-beam search.

\subsection{Main Results: Physical Grounding and CSI-to-Language}

\subsubsection{Comparison with Conventional Physical Estimators}

We compare \model{} with representative learned and analytical CSI estimators: PDP/IFFT-based and
flattened-CSI MLPs, a generic Transformer, a beam--frequency CNN, single-task CSI heads, and a
non-learned IFFT--PDP estimator. Together, these baselines test the effects of delay-domain processing,
shared representation learning, and target-specific prediction branches.

All methods are evaluated on the D2Los\_Data split under the same UPA-$8{\times}8$ setting with
128 CSI frequency points. For each seed, 40,000 training samples are randomly drawn from the source
pool and evaluated on 10,000 held-out samples. LoS and NLoS links are balanced in both training and
test sets to avoid visibility bias, and learned models are trained with three random seeds. As shown
in Figure~\ref{fig:method-overview}, \model{} uses the same CSI observation as input but routes the
shared representation through physics-aware branches before deterministic report generation.

\begin{table}[!t]
\centering
\caption{Main balanced-split physical results. Learned methods report mean $\pm$ sample standard
deviation over three seeds; the deterministic IFFT--PDP baseline is evaluated once on 10,000 test
examples. Lower is better for MAE; higher is better for reflection-bin accuracy. Underlining marks
the best result across all methods, including ours. The final row gives the relative change from the
best prior method, with direction indicated by color (red: improvement; blue: degradation).}
\label{tab:physical-main}
\scriptsize
\setlength{\tabcolsep}{4.0pt}
\resizebox{\columnwidth}{!}{%
\begin{tabular}{lccc}
\toprule
\textbf{Method} & \textbf{Delay (ns) $\downarrow$} & \textbf{Angle ($^\circ$) $\downarrow$} &
\textbf{Power (dB) $\downarrow$} \\
\midrule
\rowcolor{modelblue}\textbf{\model{}} & $\best{\boldsymbol{21.911\!\pm\!0.303}}$ & $\best{\boldsymbol{18.372\!\pm\!0.365}}$ & $\best{\boldsymbol{4.958\!\pm\!0.269}}$ \\
\textbf{CSI encoder + single-task heads} & $51.164\!\pm\!6.248$ & $20.449\!\pm\!0.431$ & $5.074\!\pm\!0.209$ \\
\textbf{PDP/IFFT + MLP} & $28.750\!\pm\!4.611$ & $40.292\!\pm\!0.159$ & $12.236\!\pm\!0.278$ \\
\textbf{Flattened CSI + MLP} & $114.371\!\pm\!0.978$ & $32.782\!\pm\!0.413$ & $13.315\!\pm\!1.221$ \\
\textbf{Transformer without branches} & $145.177\!\pm\!0.003$ & $42.614\!\pm\!0.005$ & $28.091\!\pm\!0.012$ \\
\textbf{CNN baseline} & $169.671\!\pm\!0.021$ & $53.111\!\pm\!0.005$ & $38.090\!\pm\!0.013$ \\
\textbf{IFFT--PDP (non-learned)} & $160.262$ & N/A & $6.555$ \\
\midrule
\rowcolor{gainrow}\textbf{Rel. change vs. best prior} & \relgain{$23.8\%$} & \relgain{$10.2\%$} & \relgain{$2.3\%$} \\
\bottomrule
\end{tabular}}
\vspace{2pt}

\resizebox{\columnwidth}{!}{%
\begin{tabular}{lccc}
\toprule
\textbf{Method} & \textbf{$K$ (dB) $\downarrow$} & \textbf{Refl. MAE $\downarrow$} &
\textbf{Bin acc. $\uparrow$} \\
\midrule
\rowcolor{modelblue}\textbf{\model{}} & $\boldsymbol{1.983\!\pm\!0.263}$ & $\best{\boldsymbol{0.983\!\pm\!0.005}}$ & $\best{\boldsymbol{0.7061\!\pm\!0.0039}}$ \\
\textbf{CSI encoder + single-task heads} & $\best{1.521\!\pm\!0.090}$ & $1.222\!\pm\!0.072$ & $0.5523\!\pm\!0.0259$ \\
\textbf{PDP/IFFT + MLP} & $5.002\!\pm\!0.275$ & $1.706\!\pm\!0.032$ & $0.4420\!\pm\!0.0201$ \\
\textbf{Flattened CSI + MLP} & $5.600\!\pm\!1.013$ & $1.835\!\pm\!0.196$ & $0.4237\!\pm\!0.0999$ \\
\textbf{Transformer without branches} & $20.762\!\pm\!0.051$ & $3.214\!\pm\!0.021$ & $0.1808\!\pm\!0.0008$ \\
\textbf{CNN baseline} & $40.852\!\pm\!0.055$ & $10.134\!\pm\!0.011$ & $0.1562\!\pm\!0.1463$ \\
\textbf{IFFT--PDP (non-learned)} & $20.063^{\ast}$ & N/A & N/A \\
\midrule
\rowcolor{gainrow}\textbf{Rel. change vs. best prior} & \relloss{$30.4\%$} & \relgain{$19.6\%$} & \relgain{$27.8\%$} \\
\bottomrule
\end{tabular}}
\vspace{2pt}
\end{table}

Table~\ref{tab:physical-main} shows that conventional estimators can fit individual attributes but
struggle to provide a complete and consistent propagation record. The non-learned IFFT--PDP pipeline
is useful for direct delay-domain cues, yet it cannot output receiver-side angle or reflection fields.
Flattened MLP and CNN models are much weaker on delay, angle, power, and multipath quantities because
they do not explicitly preserve the coupled beam--frequency evidence. The generic Transformer has
higher capacity but still lacks the target-matched temporal, angular, power, and structure branches.
By contrast, \model{} obtains the best delay, angle, power, reflection-MAE, and reflection-bin
accuracy, with relative gains of 23.8\%, 10.2\%, 2.3\%, 19.6\%, and 27.8\%, respectively, over the
best prior baseline. The only exception is the scalar $K$-factor result of the single-task baseline,
which indicates that the proposed design is strongest for coupled channel reporting rather than for
every isolated scalar target.

\subsubsection{Comparison with General-Purpose Language Models}

We next compare with direct CSI-to-language baselines built on general-purpose LLMs. Qwen3-1.7B and
Qwen3.5-2B are compact open language models, InternLM3-8B provides a larger instruction-following
backbone, and DeepSeek-R1-Qwen3-8B represents a reasoning-oriented decoder. These models are tested
as direct report generators rather than as post-processing modules for the structured predictions of
\model{}.

The direct-decoding route maps CSI encoder features through an MLP into continuous prefix embeddings
and asks each LLM to produce a fixed-schema physical JSON record and textual description. Training
supervises only the target response, and the comparison uses the same balanced protocol and common
1,000-example cohort as \model{}.

\begin{table}[H]
\centering
\caption{Structured prediction versus direct general-purpose language decoding on the common
1,000-example cohort. Panel A reports description-level quality; Panel B reports field-level error.
Higher is better for accuracy and slot F1; lower is better for hallucination and MAE. Underlining
marks the best result across all methods, including ours. Signed changes are computed against the
best direct decoder; direction is indicated by color (red: improvement; blue: degradation).}
\label{tab:decoder-overall}
\scriptsize
\setlength{\tabcolsep}{4pt}
{\scriptsize\textbf{Panel A: description-level quality}}\par\vspace{1pt}
\resizebox{\columnwidth}{!}{%
\begin{tabular}{lccc}
\toprule
\textbf{Method} & \textbf{Factuality (\%) $\uparrow$} & \textbf{Numeric acc. (\%) $\uparrow$} &
\textbf{Slot F1 (\%) $\uparrow$} \\
\midrule
\rowcolor{modelblue}\textbf{\model{}} & \best{\textbf{96.43}} & \best{\textbf{82.61}} & \best{\textbf{99.96}} \\
\textbf{Qwen3-1.7B} & 92.47 & 64.15 & 99.85 \\
\textbf{Qwen3.5-2B} & 92.68 & 65.06 & 99.86 \\
\textbf{InternLM3-8B} & 92.76 & 65.71 & 99.84 \\
\textbf{DeepSeek-R1-Qwen3-8B} & 83.27 & 44.72 & 97.82 \\
\midrule
\rowcolor{gainrow}\textbf{Rel. change vs. best prior} & \relgain{$4.0\%$} & \relgain{$25.7\%$} & \relgain{$0.1\%$} \\
\bottomrule
\end{tabular}}
\vspace{3pt}

\resizebox{\columnwidth}{!}{%
\begin{tabular}{lccc}
\toprule
\textbf{Method} & \textbf{Halluc. (\%) $\downarrow$} &
\textbf{Delay acc. (\%) $\uparrow$} & \textbf{Delay MAE (ns) $\downarrow$} \\
\midrule
\rowcolor{modelblue}\textbf{\model{}} & \best{\textbf{0.074}} & \best{\textbf{96.05}} & \best{\textbf{21.91}} \\
\textbf{Qwen3-1.7B} & 0.14 & 22.60 & 177.66 \\
\textbf{Qwen3.5-2B} & 0.14 & 28.33 & 147.39 \\
\textbf{InternLM3-8B} & 0.16 & 26.45 & 165.56 \\
\textbf{DeepSeek-R1-Qwen3-8B} & 4.00 & 17.81 & 232.66 \\
\midrule
\rowcolor{gainrow}\textbf{Rel. change vs. best prior} & \relgain{$47.1\%$} & \relgain{$239.0\%$} & \relgain{$85.1\%$} \\
\bottomrule
\end{tabular}}
\vspace{3pt}

{\scriptsize\textbf{Panel B: field-level MAE}}\par\vspace{1pt}
\setlength{\tabcolsep}{2.2pt}
\resizebox{\columnwidth}{!}{%
\begin{tabular}{lccccc}
\toprule
\textbf{Method} & \textbf{\shortstack{First delay\\(ns) $\downarrow$}} &
\textbf{\shortstack{\los{} delay\\(ns) $\downarrow$}} &
\textbf{\shortstack{Spread\\(ns) $\downarrow$}} &
\textbf{\shortstack{First angle\\($^\circ$) $\downarrow$}} &
\textbf{\shortstack{\los{} angle\\($^\circ$) $\downarrow$}} \\
\midrule
\rowcolor{modelblue}\textbf{\model{}} & $\best{\boldsymbol{24.75\!\pm\!1.04}}$ & $\best{\boldsymbol{8.83\!\pm\!0.55}}$ & $\best{\boldsymbol{11.30\!\pm\!0.26}}$ & $\best{\boldsymbol{16.61\!\pm\!0.56}}$ & $\best{\boldsymbol{1.33\!\pm\!0.17}}$ \\
Qwen3-1.7B & 177.66 & 146.54 & 47.03 & 42.84 & 28.71 \\
Qwen3.5-2B & 147.39 & 89.09 & 44.48 & 40.60 & 27.66 \\
InternLM3-8B & 165.56 & 127.86 & 46.57 & 40.22 & 29.07 \\
DeepSeek-R1-Qwen3-8B & 232.66 & 175.22 & 49.64 & 59.16 & 36.05 \\
\midrule
\rowcolor{gray!8}\textbf{Rel. change vs. best direct LLM} & \relgain{$83.2\%$} & \relgain{$90.1\%$} & \relgain{$74.6\%$} & \relgain{$58.7\%$} & \relgain{$95.2\%$} \\
\bottomrule
\end{tabular}}

\vspace{3pt}
\resizebox{\columnwidth}{!}{%
\begin{tabular}{lccccc}
\toprule
\textbf{Method} & \textbf{\shortstack{Az. spread\\($^\circ$) $\downarrow$}} &
\textbf{\shortstack{First power\\(dB) $\downarrow$}} &
\textbf{\shortstack{$K$-factor\\(dB) $\downarrow$}} &
\textbf{Refl. MAE $\downarrow$} & \textbf{Path MAE $\downarrow$} \\
\midrule
\rowcolor{modelblue}\textbf{\model{}} & $\boldsymbol{9.72\!\pm\!0.47}$ & $\best{\boldsymbol{4.34\!\pm\!0.27}}$ & $\best{\boldsymbol{2.15\!\pm\!0.04}}$ & $\best{\boldsymbol{1.032\!\pm\!0.019}}$ & $\best{\boldsymbol{0.086\!\pm\!0.005}}$ \\
Qwen3-1.7B & 8.16 & 6.17 & 2.559 & 1.409 & 0.157 \\
Qwen3.5-2B & \best{6.76} & 5.80 & 2.360 & 1.443 & 0.160 \\
InternLM3-8B & 8.10 & 5.96 & 2.485 & 1.494 & 0.163 \\
DeepSeek-R1-Qwen3-8B & 53.52 & 27.77 & 16.418 & 3.618 & 0.578 \\
\midrule
\rowcolor{gray!8}\textbf{Rel. change vs. best direct LLM} & \relloss{$43.8\%$} & \relgain{$25.2\%$} & \relgain{$8.9\%$} & \relgain{$26.8\%$} & \relgain{$45.2\%$} \\
\bottomrule
\end{tabular}}
\label{tab:decoder-fields}
\end{table}

Table~\ref{tab:decoder-overall} shows that direct LLM decoders can learn the output schema and produce
fluent-looking channel reports, but their numerical grounding is much weaker. The best direct decoder
reaches only 65.71\% numerical accuracy, whereas \model{} reaches 82.61\%. The gap is more pronounced
for physical fields: \model{} reduces first-delay MAE by 83.2\%, LoS-delay MAE by 90.1\%, first-angle
MAE by 58.7\%, and LoS-angle MAE by 95.2\% relative to the best direct LLM. It also lowers first-power,
$K$-factor, reflection, and path-count errors. These results indicate that the main advantage of
\model{} is not text fluency alone, but the separation of physical inference, relation correction, and
deterministic verbalization.

\subsubsection{Measured Inference Cost}

A deployable channel-reporting system must be accurate without imposing the latency and memory
footprint of autoregressive generation. We therefore conduct a controlled inference-cost measurement:
all methods are evaluated on the same CSI test samples, the same GPU, and a fixed warmup/repeat
protocol. We record end-to-end latency, CUDA time, throughput, peak allocated VRAM, parameter count,
and generated-token statistics, and then merge these measurements with the corresponding physical and
text-quality metrics. We therefore measure end-to-end latency on 200 common examples and peak
allocated VRAM under explicit precision settings; the quality results use the separate 1,000-example
decoder cohort.

\begin{table}[H]
\centering
\caption{Measured end-to-end inference cost. Conventional models use FP32, whereas direct language
decoders use 4-bit NF4 weights with BF16 computation. Lower is better. Underlining marks the best
result across all methods, including ours. The final row compares \model{} with the most efficient
direct language decoder for each cost measure.}
\label{tab:inference-cost-main}
\scriptsize
\setlength{\tabcolsep}{2.2pt}
\resizebox{\columnwidth}{!}{%
\begin{tabular}{lcccc}
\toprule
\textbf{Method} & \textbf{Median (ms) $\downarrow$} & \textbf{P95 (ms) $\downarrow$} &
\textbf{VRAM (GB) $\downarrow$} & \textbf{Params $\downarrow$} \\
\midrule
\rowcolor{modelblue}\textbf{CSI2TEXT-PGMT-COST-BENCHMARK} & \textbf{13.58} & \textbf{14.41} & \textbf{0.13} & \textbf{32.1M} \\
\textbf{Single-task} & 13.47 & 13.78 & 0.23 & 58.9M \\
\textbf{PDP/IFFT+MLP} & 3.76 & 3.85 & \best{0.02} & \best{3.3M} \\
\textbf{Flattened MLP} & \best{1.63} & \best{1.82} & 0.33 & 86.5M \\
\textbf{Transformer} & 7.44 & 7.58 & 0.11 & 28.1M \\
\textbf{CNN} & 3.44 & 3.49 & 0.04 & 6.5M \\
\textbf{Qwen3-1.7B} & 12446 & 15391 & 1.53 & 1.07B \\
\textbf{Qwen3.5-2B} & 11656 & 14450 & 2.03 & 1.41B \\
\textbf{DeepSeek-8B} & 16905 & 33112 & 6.18 & 4.81B \\
\textbf{InternLM3-8B} & 23046 & 27953 & 6.18 & 5.03B \\
\midrule
\rowcolor{gainrow}\textbf{Rel. change vs. best direct LLM} & \relgain{$99.9\%$} & \relgain{$99.9\%$} & \relgain{$91.5\%$} & \relgain{$97.0\%$} \\
\bottomrule
\end{tabular}}
\end{table}

\begin{figure}[H]
\centering
\caption{Quality--cost Pareto comparison. The left panel relates description factuality to median
inference latency, and the right panel relates numerical tolerance accuracy to latency. Bubble area
encodes peak allocated VRAM. The structured model occupies the high-quality, low-latency corner in
both views.}
\label{fig:quality-cost-pareto}
\includegraphics[width=\columnwidth]{quality_cost_pareto_horizontal.pdf}
\end{figure}

Table~\ref{tab:inference-cost-main} and Figure~\ref{fig:quality-cost-pareto} show that direct language
decoding pays a large systems cost in addition to its weaker numerical grounding. The direct decoders
operate on a seconds-per-report scale and allocate far more memory, whereas \model{} remains in the
millisecond regime with a compact physical predictor. Conventional regressors can be slightly faster,
but they do not approach the complete quality profile in Table~\ref{tab:physical-main}. The proposed
pipeline therefore occupies the favorable accuracy--cost corner: it preserves structured factuality
without requiring autoregressive generation for every report.

This distinction matters for radio-network operation, where many candidate links or beams may need
to be evaluated within one scheduling interval. The LLM baselines must keep large backbones resident
and generate output tokens even for a fixed report schema. In contrast, \model{} predicts the
structured record in one feed-forward pass and then uses deterministic verbalization, retaining
physically auditable reports while staying close to the cost profile of compact neural estimators.

\subsection{Generalization}

\subsubsection{Frequency-Configuration Generalization}

\begin{figure}[H]
\centering
\makebox[\columnwidth][c]{\includegraphics[width=1.00\columnwidth]{figures/figure4.png}}
\caption{Leave-one-$N_f$-out generalization. The single-column two-panel line plot reports
performance degradation relative to the seen-$N_f=128$ reference across held-out frequency
samplings. Panel (a) summarizes accuracy, factuality, numeric accuracy, and $K$-factor trends,
whereas panel (b) highlights first-path angle, delay, and power. Positive values indicate worse
held-out performance.}
\label{fig:frequency-generalization}
\end{figure}

For each column of Figure~\ref{fig:frequency-generalization}, the complete model is trained on the
other four values in $\{64,96,128,192,256\}$ and then evaluated on the excluded sampling
configuration. Visibility classification and report factuality remain stable across the held-out
settings, while fine temporal localization is the most sensitive quantity and power or angle exhibit
smaller configuration-dependent changes. The result shows that the configuration-aware
representation transfers semantic channel structure more reliably than it transfers exact delay
resolution. An assisted setting with paired observations from two numerologies is examined separately
as a complementary multi-view extension.

\subsubsection{Cross-Array and Directional Transfer}

We train a mixed-array model with UPA-$8{\times}8$ and ULA-64 observations, then evaluate it on
those seen geometries and on aligned, held-out UPA-$4{\times}4$, UPA-$16{\times}4$, and ULA-32
groups without geometry-specific fine-tuning. A separate pair of experiments reverses the ULA
transfer direction, comparing larger-to-smaller and smaller-to-larger apertures while keeping the
evaluation groups disjoint from training.

\begin{table}[H]
\centering
\caption{Cross-array and directional ULA transfer over three seeds. UPA-$8{\times}8$ and ULA-64 are
seen during mixed-array training. Lower is better for MAE; higher is better for factuality. Panel B
also reports the direction-aware relative change from the seen setting to the unseen setting.}
\label{tab:cross-array}
\scriptsize
\setlength{\tabcolsep}{2.7pt}
\textbf{Panel A: cross-array transfer}\par\vspace{2pt}
\resizebox{\columnwidth}{!}{%
\begin{tabular}{lccccc}
\toprule
\textbf{Array} & \textbf{Status} & \textbf{LoS angle $\downarrow$} &
\textbf{NLoS angle $\downarrow$} & \textbf{First delay $\downarrow$} &
\textbf{Factuality $\uparrow$} \\
\midrule
UPA-$8{\times}8$ & Seen & $6.106\!\pm\!1.966$ & $42.591\!\pm\!0.377$ & $53.854\!\pm\!6.213$ & $94.316\!\pm\!0.096\%$ \\
ULA-64 & Seen & $14.566\!\pm\!2.623$ & $42.710\!\pm\!0.562$ & $47.449\!\pm\!2.121$ & $94.178\!\pm\!0.167\%$ \\
UPA-$4{\times}4$ & Unseen & $10.257\!\pm\!3.013$ & $45.859\!\pm\!0.110$ & $44.082\!\pm\!0.786$ & $93.556\!\pm\!0.092\%$ \\
UPA-$16{\times}4$ & Unseen & $9.446\!\pm\!0.913$ & $43.237\!\pm\!0.114$ & $54.828\!\pm\!5.982$ & $93.639\!\pm\!0.136\%$ \\
ULA-32 & Unseen & $11.345\!\pm\!0.566$ & $42.919\!\pm\!0.400$ & $43.399\!\pm\!1.029$ & $94.449\!\pm\!0.193\%$ \\
\bottomrule
\end{tabular}}
\vspace{3pt}

\textbf{Panel B: directional ULA transfer}\par\vspace{2pt}
\resizebox{\columnwidth}{!}{%
\begin{tabular}{lccccc}
\toprule
\textbf{Direction} & \textbf{Status} & \textbf{LoS angle $\downarrow$} &
\textbf{NLoS angle $\downarrow$} & \textbf{First delay $\downarrow$} &
\textbf{Factuality $\uparrow$} \\
\midrule
$64\rightarrow32$ & Seen 64 & $14.566\!\pm\!2.623$ & $45.613\!\pm\!0.322$ & $47.449\!\pm\!2.121$ & $94.178\!\pm\!0.167\%$ \\
 & Unseen 32 & $11.345\!\pm\!0.566$ & $41.740\!\pm\!0.462$ & $43.399\!\pm\!1.029$ & $94.449\!\pm\!0.193\%$ \\
\rowcolor{gainrow} & \textbf{Unseen vs. seen} & \relgain{22.1\%} & \relgain{8.5\%} & \relgain{8.5\%} & \relgain{0.3\%} \\
\midrule
$32\rightarrow64$ & Seen 32 & $10.730\!\pm\!0.785$ & $42.982\!\pm\!0.112$ & $41.619\!\pm\!0.954$ & $95.092\!\pm\!0.096\%$ \\
 & Unseen 64 & $17.223\!\pm\!3.331$ & $48.592\!\pm\!0.628$ & $49.350\!\pm\!3.524$ & $94.222\!\pm\!0.073\%$ \\
\rowcolor{gainrow} & \textbf{Unseen vs. seen} & \relloss{60.5\%} & \relloss{13.1\%} & \relloss{18.6\%} & \relloss{0.9\%} \\
\bottomrule
\end{tabular}}
\label{tab:directional-ula}
\end{table}

Table~\ref{tab:cross-array} shows that factual report generation remains stable on unseen array shapes
and sizes, while the persistent NLoS-angle error is shared across geometries rather than caused by one
specific array. Directional transfer is asymmetric: learning from the larger aperture transfers cleanly
to the smaller one, whereas extrapolating from a smaller to a larger aperture produces a clearer loss
in angular and temporal precision. The model therefore generalizes across array configurations, but
the training aperture still bounds the reliability of geometric extrapolation.

\subsubsection{Spatial-Block and Map-Disjoint Generalization}

To distinguish local interpolation from environment transfer, we compare a random split with
$32\,\mathrm{m}\times32\,\mathrm{m}$ spatial blocks and complete held-out maps. The audited cohorts
have zero overlap in maps, map-qualified blocks, and user-equipment positions, so the stricter
protocols cannot reuse nearby or identical propagation links from training.

\begin{table}[H]
\centering
\caption{Random, spatial-block, and complete-map-disjoint evaluation. The final column reports the
direction-aware performance difference between map-disjoint evaluation and random splitting.
Higher is better for accuracy; lower is better for MAE. Red denotes an improvement and blue denotes
a degradation. Underlining marks the best result in each row, and boldface highlights the
map-disjoint results.}
\label{tab:split-sensitivity}
\scriptsize
\setlength{\tabcolsep}{2.5pt}
\resizebox{\columnwidth}{!}{%
\begin{tabular}{lcccc}
\toprule
Metric & Random & Spatial block & Map-disjoint & Map-disjoint vs. Random \\
\midrule
Semantic accuracy $\uparrow$ & 0.8750 & \underline{0.8763} & \textbf{0.8618} & \relloss{$1.32$ pp} \\
Semantic macro-F1 $\uparrow$ & 0.5960 & \underline{0.6352} & \textbf{0.4884} & \relloss{$10.76$ pp} \\
First delay (ns) $\downarrow$ & \underline{25.8400} & 31.7410 & \textbf{38.5178} & \relloss{$49.1\%$} \\
LoS delay (ns) $\downarrow$ & \underline{10.2439} & 13.4182 & \textbf{17.5738} & \relloss{$71.6\%$} \\
NLoS delay (ns) $\downarrow$ & \underline{47.4350} & 50.0639 & \textbf{59.4618} & \relloss{$25.4\%$} \\
Delay spread (ns) $\downarrow$ & \underline{14.9583} & 16.5812 & \textbf{18.6648} & \relloss{$24.8\%$} \\
$K$ (dB) $\downarrow$ & \underline{2.4280} & 2.7301 & \textbf{3.2529} & \relloss{$34.0\%$} \\
First power (dB) $\downarrow$ & \underline{4.6009} & 5.2943 & \textbf{6.2998} & \relloss{$36.9\%$} \\
First angle ($^\circ$) $\downarrow$ & 21.8495 & 23.9274 & \textbf{\underline{21.7308}} & \relgain{$0.5\%$} \\
Azimuth spread ($^\circ$) $\downarrow$ & \underline{7.5784} & 10.1609 & \textbf{10.0201} & \relloss{$32.2\%$} \\
Reflection-bin accuracy $\uparrow$ & \underline{0.6599} & 0.5711 & \textbf{0.5498} & \relloss{$11.01$ pp} \\
Path-count MAE $\downarrow$ & 0.1965 & 0.2321 & \textbf{\underline{0.1878}} & \relgain{$4.4\%$} \\
\bottomrule
\end{tabular}}
\end{table}

Table~\ref{tab:split-sensitivity} shows that performance remains usable when evaluation moves from
random interpolation to spatially or map-disjoint transfer. Delay, power, dispersion, and reflection
are more sensitive to unseen environments, whereas first-angle and path-count prediction remain
comparatively stable. The nonuniform change is important: random splitting is optimistic for selected
physical quantities, but the retained semantic and structural performance demonstrates meaningful
generalization beyond memorized locations.

\subsection{Robustness}

We evaluate robustness by perturbing the observed complex CSI with additive white Gaussian noise
(AWGN) during both training and testing. For each sample, noise is applied only to the valid,
non-padding index set \(\Omega\). Given a target signal-to-noise ratio \(\gamma\) in decibels, the
corrupted CSI is
\begin{equation}
\begin{aligned}
\widetilde{\mathbf H}&=\mathbf H+\mathbf N, &
P_H&=\frac{1}{|\Omega|}\sum_{u\in\Omega}|H_u|^2,\\
P_N&=P_H10^{-\gamma/10}, &
\Re(N_u),\Im(N_u)&\sim\mathcal N(0,P_N/2),
\end{aligned}
\label{eq:complex-awgn}
\end{equation}
where \(\mathbb E|N_u|^2=P_N\). During noise-augmented training, each batch is independently
selected for corruption with \(z_b\sim\operatorname{Bernoulli}(0.8)\), and its SNR is sampled as
\(\gamma_b\sim\mathcal U(10,30)\) dB. After warm-up, the batch objective is
\begin{equation}
\begin{aligned}
\mathcal L_b={}&\mathcal L_{\mathrm{clean}}+0.2\mathcal L_{K,\mathrm{clean}}\\
&+z_b\bigl[0.5\mathcal L_{\mathrm{noisy}}+0.1\mathcal L_{\tau,\mathrm{cons}}\\
&\hspace{18mm}+0.1\mathcal L_{K,\mathrm{cons}}
+0.1\mathcal L_{K,\mathrm{noisy}}\bigr],
\end{aligned}
\label{eq:noise-training}
\end{equation}
where the supervised terms are masked Smooth-L1 losses and the consistency terms compare noisy
predictions with stop-gradient clean predictions in normalized physical-target space. Testing uses
freshly sampled noise at \(\gamma\in\{30,20,10,0\}\) dB, together with an uncorrupted reference set,
so the evaluation measures generalization to unseen noise realizations rather than reuse of the
training perturbations.

For each trained model seed \(s\), every metric is first averaged over three independent test-noise
realizations \(r\), and the final mean and sample standard deviation are then computed over three
model seeds:
\begin{equation}
\begin{aligned}
\bar M_s&=\frac{1}{3}\sum_{r=1}^{3}M_{s,r},&
\mu_M&=\frac{1}{3}\sum_{s=1}^{3}\bar M_s,\\
\sigma_M&=\sqrt{\frac{1}{2}\sum_{s=1}^{3}(\bar M_s-\mu_M)^2}.
\end{aligned}
\label{eq:awgn-seed-statistics}
\end{equation}
For a lower-is-better metric, the relative degradation at SNR \(\gamma\) is
\begin{equation}
D_M(\gamma)=
\frac{M(\gamma)-M(\mathrm{clean})}{M(\mathrm{clean})}\times100\%.
\label{eq:relative-noise-degradation}
\end{equation}

\begin{table}[H]
\centering
\caption{Three-seed test-time complex-AWGN robustness. Lower is better for MAE; higher is better
for numerical tolerance accuracy.}
\label{tab:awgn}
\scriptsize
\setlength{\tabcolsep}{4pt}
\resizebox{\columnwidth}{!}{%
\begin{tabular}{lccccc}
\toprule
Metric & Clean & 30 dB & 20 dB & 10 dB & 0 dB \\
\midrule
First delay (ns) $\downarrow$ & $25.97\!\pm\!0.34$ & $72.70\!\pm\!0.40$ & $86.11\!\pm\!0.96$ & $100.23\!\pm\!1.73$ & $209.29\!\pm\!12.31$ \\
LoS delay (ns) $\downarrow$ & $8.72\!\pm\!0.49$ & $20.30\!\pm\!2.71$ & $25.00\!\pm\!2.43$ & $34.10\!\pm\!1.32$ & $231.85\!\pm\!28.52$ \\
NLoS delay (ns) $\downarrow$ & $43.22\!\pm\!0.27$ & $125.10\!\pm\!3.05$ & $147.22\!\pm\!2.52$ & $166.36\!\pm\!4.31$ & $186.73\!\pm\!6.54$ \\
First angle ($^\circ$) $\downarrow$ & $18.84\!\pm\!0.40$ & $23.10\!\pm\!0.84$ & $24.29\!\pm\!0.71$ & $25.12\!\pm\!0.79$ & $32.30\!\pm\!1.92$ \\
First power (dB) $\downarrow$ & $5.12\!\pm\!0.35$ & $4.88\!\pm\!0.37$ & $4.79\!\pm\!0.37$ & $5.38\!\pm\!0.58$ & $17.69\!\pm\!0.08$ \\
$K$ (dB) $\downarrow$ & $2.084\!\pm\!0.124$ & $2.142\!\pm\!0.125$ & $2.612\!\pm\!0.118$ & $5.978\!\pm\!0.582$ & $19.855\!\pm\!0.659$ \\
Reflection-path MAE $\downarrow$ & $1.018\!\pm\!0.014$ & $1.180\!\pm\!0.006$ & $1.232\!\pm\!0.011$ & $1.338\!\pm\!0.009$ & $2.384\!\pm\!0.178$ \\
Numeric accuracy (\%) $\uparrow$ & $74.05\!\pm\!0.57$ & $69.27\!\pm\!0.65$ & $64.91\!\pm\!0.82$ & $56.22\!\pm\!1.01$ & $33.93\!\pm\!2.18$ \\
\bottomrule
\end{tabular}}
\end{table}

\begin{figure}[H]
\centering
\caption{Test-time complex-AWGN robustness. Panel (a) reports the absolute MAE of first-path,
LoS, and NLoS delay; panel (b) reports first-angle, first-power, and \(K\)-factor MAE normalized by
their clean-condition values, focusing the comparison on physical prediction robustness.}
\label{fig:awgn-robustness}
\makebox[\columnwidth][c]{\includegraphics[width=1.00\columnwidth]{figures/awgn_robustness_replacement.png}}
\end{figure}

Table~\ref{tab:awgn} and Figure~\ref{fig:awgn-robustness} show that degradation is monotonic overall
but strongly attribute dependent. Delay is affected first because it is encoded by phase variation
across frequency. The NLoS delay curve rises much faster than the LoS curve at moderate SNR, which
is consistent with the difficulty of locating the earliest indirect component among several
overlapping reflections. First-angle error grows more gradually, indicating that coarse spatial
structure survives moderate corruption better than fine temporal localization.

Power and \(K\)-factor remain close to their clean baselines at moderate noise levels but deteriorate
sharply when the SNR approaches 0 dB. The small fluctuation in power error at 20--30 dB lies within
the reported seed variation and should not be interpreted as an improvement. Reflection-path error
also increases steadily, while numerical tolerance accuracy in Table~\ref{tab:awgn} decreases as more
fields leave their accepted ranges. Overall, noise augmentation preserves useful physical predictions
under moderate corruption, but it cannot recover channel evidence destroyed by severe noise. This
attribute-specific behavior supports confidence-aware use: coarse state may remain usable even when
a precise delay or power value should be withheld.

\subsection{Ablation}

\subsubsection{Language-Alignment Ablation}

We construct matched variants that remove CSI--text alignment, text--prototype alignment, or both
(the physics-only setting), while keeping the physical heads, target supervision, correction rules,
data split, and verbalizer fixed. The comparison therefore isolates training-time semantic alignment
without changing the physical predictor or report-generation protocol.

\begin{table}[H]
\centering
\caption{Matched language-alignment ablations over three seeds. Higher is better for classification
and bin accuracy; lower is better for MAE. The final row in each block compares the full model with
the best ablated variant for each metric (red: improvement; blue: degradation). Underlining marks
the best result across all variants, including the full model.}
\label{tab:language-ablation}
\scriptsize
\setlength{\tabcolsep}{2.7pt}
\resizebox{\columnwidth}{!}{%
\begin{tabular}{lcccc}
\toprule
\textbf{Variant} & \textbf{LoS/NLoS acc. $\uparrow$} & \textbf{Delay (ns) $\downarrow$} &
\textbf{Refl. MAE $\downarrow$} & \textbf{Refl. acc. $\uparrow$} \\
\midrule
\rowcolor{modelblue}\textbf{Full alignment} & $\best{\boldsymbol{0.99572\!\pm\!0.00314}}$ & $\best{\boldsymbol{25.9495\!\pm\!1.2610}}$ & $\best{\boldsymbol{1.0644\!\pm\!0.0164}}$ & $\best{\boldsymbol{0.6866\!\pm\!0.0053}}$ \\
\textbf{No CSI--text} & $0.99297\!\pm\!0.00080$ & $27.3459\!\pm\!0.9163$ & $1.0776\!\pm\!0.0069$ & $0.6579\!\pm\!0.0022$ \\
\textbf{No text--prototype} & $0.99330\!\pm\!0.00006$ & $28.3746\!\pm\!0.2799$ & $1.0925\!\pm\!0.0100$ & $0.6521\!\pm\!0.0092$ \\
\textbf{Physics only} & $0.98313\!\pm\!0.00071$ & $27.6446\!\pm\!0.8540$ & $1.1170\!\pm\!0.0135$ & $0.6280\!\pm\!0.0062$ \\
\midrule
\rowcolor{gainrow}\textbf{Rel. change vs. best ablation} & \relgain{$0.2\%$} & \relgain{$5.1\%$} & \relgain{$1.2\%$} & \relgain{$4.4\%$} \\
\bottomrule
\end{tabular}}
\vspace{2pt}

\resizebox{\columnwidth}{!}{%
\begin{tabular}{lccc}
\toprule
\textbf{Variant} & \textbf{Path MAE $\downarrow$} &
\textbf{NLoS power $\downarrow$} & \textbf{NLoS angle $\downarrow$} \\
\midrule
\rowcolor{modelblue}\textbf{Full alignment} & $\best{\boldsymbol{0.2106\!\pm\!0.0102}}$ & $\best{\boldsymbol{7.3023\!\pm\!0.2223}}$ & $\best{\boldsymbol{38.3226\!\pm\!0.7280}}$ \\
\textbf{No CSI--text} & $0.2130\!\pm\!0.0144$ & $7.4089\!\pm\!0.1944$ & $39.6833\!\pm\!0.7642$ \\
\textbf{No text--prototype} & $0.2144\!\pm\!0.0559$ & $7.3636\!\pm\!0.0424$ & $39.6135\!\pm\!1.2067$ \\
\textbf{Physics only} & $0.3736\!\pm\!0.0551$ & $8.3348\!\pm\!0.5728$ & $39.2445\!\pm\!1.3988$ \\
\midrule
\rowcolor{gainrow}\textbf{Rel. change vs. best ablation} & \relgain{$1.1\%$} & \relgain{$0.8\%$} & \relgain{$2.3\%$} \\
\bottomrule
\end{tabular}}
\end{table}

Full alignment gives the strongest overall results, with the clearest gains in LoS/NLoS
classification, reflection-bin accuracy, and path-count estimation. Frozen canonical text anchors
preserve the predefined semantic reference, while adaptive prototypes adjust class centers to the
observed CSI distribution. Their complementary effect supports semantic alignment as a structural
regularizer rather than a replacement for direct physical supervision.

\subsubsection{Physical Prediction-Head Ablations}

The physical estimator is deliberately factorized into a shared configuration encoder and shared
physics token, followed by multiple target-specific encoders and prediction heads for delay, angle,
power, and path structure. This multi-encoder, multi-head design allows coupled evidence to be reused
while retaining the inductive bias appropriate to each physical quantity. We remove one component or
auxiliary constraint at a time and evaluate only the labels most directly related to that intervention.
Table~\ref{tab:physical-head-ablation} combines the angle, structural, and temporal/power studies.

\begin{table}[H]
\centering
\caption{Target-specific physical prediction-head ablations over three seeds. Each panel retains only
labels related to the removed component; lower is better unless indicated otherwise. The final row
compares the full model with the best ablated variant for each metric (red: improvement; blue:
degradation). Underlining marks the best result across all variants, including the full model.}
\label{tab:physical-head-ablation}
\label{tab:angle-ablation}
\label{tab:reflection-ablation}
\label{tab:architecture-ablation}
\scriptsize
\setlength{\tabcolsep}{3.0pt}
\renewcommand{\arraystretch}{0.96}
\textbf{Panel A: angle encoders}\par\vspace{1pt}
\resizebox{\linewidth}{!}{%
\begin{tabular}{lccc}
\toprule
\textbf{Variant} & \textbf{LoS angle $\downarrow$} & \textbf{First-angle LoS $\downarrow$} &
\textbf{First-angle NLoS $\downarrow$} \\
\midrule
\rowcolor{modelblue}\textbf{Full} & $\best{\boldsymbol{1.3564\!\pm\!0.1419}}$ & $\best{\boldsymbol{1.6798\!\pm\!0.0434}}$ & $\best{\boldsymbol{35.2717\!\pm\!0.7477}}$ \\
\textbf{No LoS-angle encoder} & $16.3092\!\pm\!0.2847$ & $5.9300\!\pm\!0.0964$ & $39.8998\!\pm\!1.0083$ \\
\textbf{No first-angle encoder} & $1.8599\!\pm\!0.2485$ & $2.3328\!\pm\!0.0427$ & $43.5819\!\pm\!0.0074$ \\
\midrule
\rowcolor{gainrow}\textbf{Rel. change vs. best ablation} & \relgain{$27.1\%$} & \relgain{$28.0\%$} & \relgain{$11.6\%$} \\
\bottomrule
\end{tabular}}

\vspace{2pt}
\textbf{Panel B: structural objectives}\par\vspace{1pt}
\resizebox{\linewidth}{!}{%
\begin{tabular}{lccc}
\toprule
\textbf{Variant} & \textbf{Reflection MAE $\downarrow$} &
\textbf{Reflection acc. $\uparrow$} & \textbf{Path-count MAE $\downarrow$} \\
\midrule
\rowcolor{modelblue}\textbf{Full} & $\best{\boldsymbol{0.9833\!\pm\!0.0054}}$ & $\best{\boldsymbol{0.7061\!\pm\!0.0039}}$ & $\best{\boldsymbol{0.1539\!\pm\!0.0079}}$ \\
\textbf{No reflection auxiliary loss} & $8.2848\!\pm\!0.0138$ & $0.1741\!\pm\!0.0028$ & $0.1972\!\pm\!0.0185$ \\
\textbf{No LoS consistency} & $1.0765\!\pm\!0.1479$ & $0.6548\!\pm\!0.4725$ & $0.2228\!\pm\!0.0088$ \\
\midrule
\rowcolor{gainrow}\textbf{Rel. change vs. best ablation} & \relgain{$8.7\%$} & \relgain{$7.8\%$} & \relgain{$22.0\%$} \\
\bottomrule
\end{tabular}}

\vspace{2pt}
\textbf{Panel C: shared, delay, and power components}\par\vspace{1pt}
\resizebox{\linewidth}{!}{%
\begin{tabular}{lccc}
\toprule
\textbf{Variant} & \textbf{Delay $\downarrow$} & \textbf{LoS delay $\downarrow$} &
\textbf{NLoS delay $\downarrow$} \\
\midrule
\rowcolor{modelblue}\textbf{Full} & $\best{\boldsymbol{21.9112\!\pm\!0.3031}}$ & $\best{\boldsymbol{8.1737\!\pm\!0.1556}}$ & $\best{\boldsymbol{35.8198\!\pm\!0.5986}}$ \\
\textbf{No shared token} & $22.6218\!\pm\!0.7071$ & $9.8156\!\pm\!0.6196$ & $36.6103\!\pm\!0.8507$ \\
\textbf{No delay encoder} & $69.0200\!\pm\!1.4595$ & $15.9316\!\pm\!0.8742$ & $130.8691\!\pm\!2.2043$ \\
\textbf{No power} & $29.3139\!\pm\!0.7274$ & $11.0429\!\pm\!1.0211$ & $47.7919\!\pm\!0.2816$ \\
\midrule
\rowcolor{gainrow}\textbf{Rel. change vs. best ablation} & \relgain{$3.1\%$} & \relgain{$16.7\%$} & \relgain{$2.2\%$} \\
\bottomrule
\end{tabular}}

\vspace{2pt}
\textbf{Panel D: spread, $K$, and power components}\par\vspace{1pt}
\resizebox{\linewidth}{!}{%
\begin{tabular}{lccc}
\toprule
\textbf{Variant} & \textbf{Spread $\downarrow$} & \textbf{$K$ $\downarrow$} &
\textbf{Power $\downarrow$} \\
\midrule
\rowcolor{modelblue}\textbf{Full} & $\best{\boldsymbol{11.1419\!\pm\!0.1068}}$ & $\boldsymbol{1.8829\!\pm\!0.2633}$ & $\boldsymbol{4.9575\!\pm\!0.2686}$ \\
\textbf{No shared token} & $11.1766\!\pm\!0.0753$ & $\best{1.7426\!\pm\!0.0419}$ & $\best{4.5792\!\pm\!0.0306}$ \\
\textbf{No delay encoder} & $29.6388\!\pm\!0.8355$ & $2.5399\!\pm\!0.1149$ & $5.7443\!\pm\!0.2547$ \\
\textbf{No power} & $14.1139\!\pm\!0.3321$ & $2.2289\!\pm\!0.1659$ & $4.8605\!\pm\!0.7743$ \\
\midrule
\rowcolor{gainrow}\textbf{Rel. change vs. best ablation} & \relgain{$0.3\%$} & \relloss{$8.1\%$} & \relloss{$8.3\%$} \\
\bottomrule
\end{tabular}}
\end{table}

The effects are localized to the labels matched to each removed component. Without the LoS-angle
encoder, LoS-angle estimation deteriorates sharply; without the first-angle encoder, the largest
penalty appears on NLoS first-angle prediction. Removing the reflection auxiliary objective produces
the most pronounced structural failure, simultaneously worsening reflection regression,
reflection-bin classification, and path-count estimation. The LoS-consistency term also supports
these coupled structural labels.

The delay encoder is indispensable for temporal estimation: removing it degrades first-path delay,
especially in NLoS, and also harms delay spread. Removing the raw-power branch affects not only
absolute power evidence but also delay and \(K\)-factor, confirming that scale information is shared
across the structured record. The shared token shows a smaller, mixed per-label effect, as expected
under multitask competition, but generally improves the coherence of the coupled outputs.

Taken together, the panels show why neither a single undifferentiated head nor fully independent
regressors are sufficient. Specialized heads provide target-matched representations, while shared
encoding transfers propagation-regime, path-prominence, and absolute-scale evidence across related
quantities. An occasional isolated metric improvement after removing a component does not overturn
the joint result; it reflects reduced multitask competition at the cost of other dependent labels.
The full multi-encoder, multi-head architecture is therefore necessary for a consistently accurate
and physically coherent channel record.
\subsection{Failure Modes and Error Analysis}

\begin{figure}[H]
\centering
\includegraphics[width=\columnwidth]{figures/prediction_vs_ground_truth_compact.png}
\caption{Per-sample predictions against ground truth for four principal continuous attributes. The
display contains a LoS/NLoS-stratified sample of 500 points, while the reported statistics use all
1,000 valid examples. Delay, power, and \(K\)-factor closely follow the identity line; first-angle
predictions show substantially larger NLoS dispersion. Red boxes identify representative regions
containing deviated NLoS samples. Angle error is circular.}
\label{fig:prediction-ground-truth}
\end{figure}

Figure~\ref{fig:prediction-ground-truth} first reveals that the dominant failures are concentrated in
NLoS samples rather than distributed uniformly across the test set. Delay, power, and \(K\)-factor
remain strongly ordered overall: their Pearson correlations are 0.968, 0.984, and 0.989, respectively,
and most LoS points lie close to the identity line. The largest delay outliers, the compressed
high-\(K\) region, and the broader power deviations are nevertheless predominantly NLoS cases.
First-angle prediction is the clearest failure mode. Its Pearson and Spearman correlations fall to
0.624 and 0.661, and the NLoS samples form broad off-diagonal groups even after circular angle
error is used. The figure therefore indicates that the remaining error is principally a
propagation-regime and observability problem, not a uniform calibration failure of the model.

\noindent
\begin{wrapfigure}[15]{R}{0.47\columnwidth}
\centering
\includegraphics[width=1.0\linewidth]{figures/los_nlos_ratio_replacement.png}
\captionof{figure}{Propagation‑regime error gap. Bars report
the NLoS‑to‑LoS MAE ratio; labels retain the
underlying LoS and NLoS mean $\pm$ standard deviation.}
\label{fig:los-nlos}
\end{wrapfigure}

Figure \ref{fig:los-nlos} quantifies the regime gap exposed by
Figure \ref{fig:prediction-ground-truth}. Relative to LoS, NLoS
increases first‑delay, first‑angle, and first‑power MAE by
$4.38\times$, $21.0\times$, and $2.34\times$, respectively.

Consequently, one aggregate MAE can conceal the fact that the
balanced test set combines a comparatively easy direct‑path
regime with a much harder indirect‑path regime. The particularly
large angle ratio also explains why satisfactory average
performance does not imply that every reported NLoS field is
equally reliable.

The gap mainly follows from observability. In LoS, the earliest component is tied to a direct path
with stable phase, spatial signature, and power ordering. In NLoS, the first retained component is a
reflection, and several reflected paths can have similar delays and powers, so small representation
errors may change which path is treated as dominant. The problem is amplified because the input keeps
only one receive element: receiver-side angle must be inferred indirectly from transmit-side
beam--frequency structure and statistical geometry correlations rather than resolved by a receive
array. Richer receive-side views, paired numerologies, geometry-aware priors, or NLoS-focused
training could reduce this ambiguity; when they are unavailable, high-risk NLoS cases should be
deferred or reviewed rather than forced into an automatic report, motivating the selective-prediction
mechanism below.

\subsection{Reliability-Aware Selective Prediction}

Reliable channel reporting requires more than a low average error: the system must also recognize
when an individual prediction should not be trusted. This is especially important for the NLoS tail
identified in Section~IV-E, where a small number of large delay or angle errors can affect downstream
beam management, localization, and channel interpretation even when the aggregate metric remains
acceptable. Reliability-aware selective prediction converts this requirement into an explicit
coverage--risk trade-off. Low-uncertainty reports are emitted automatically, while uncertain samples
are deferred for an additional observation, an alternative frequency configuration, or human review.
In this setting, deferral is not treated as a failure of the predictor. It is an operational decision:
the system should provide a report when the physical evidence is stable, and it should withhold or
verify a report when the model ensemble disagrees about the underlying channel attributes. This is
especially relevant for autonomous wireless controllers, where a single erroneous NLoS angle or delay
estimate can lead to a wrong beam, a poor localization update, or an overconfident channel summary.

We estimate uncertainty from disagreement among \(S=3\) independently trained models. For sample
\(i\), the ensemble prediction and its sample standard deviation are
\begin{equation}
\bar y_i=\frac{1}{S}\sum_{s=1}^{S}\hat y_i^{(s)},\qquad
u_i=\sqrt{\frac{1}{S-1}\sum_{s=1}^{S}
\left(\hat y_i^{(s)}-\bar y_i\right)^2},
\label{eq:ensemble-uncertainty}
\end{equation}
where \(u_i\) is the sample-level uncertainty score. The realized absolute error and full-coverage
MAE are
\begin{equation}
e_i=|y_i-\bar y_i|,\qquad
\operatorname{MAE}_{\mathrm{full}}=\frac{1}{N}\sum_{i=1}^{N}e_i.
\label{eq:full-coverage-error}
\end{equation}
Thus, \(\bar y_i\) is the value that would be reported if no deferral were used, and \(u_i\) measures
how sensitive that report is to the random seed and learned representation. A high \(u_i\) does not
prove that the prediction is wrong, but it indicates that the sample lies in a region where the
available CSI evidence supports several plausible physical interpretations.
After sorting the samples so that \(u_{\pi_1}\leq\cdots\leq u_{\pi_N}\), coverage \(c\) retains the
\(K(c)\) least-uncertain samples. Their selective risk is
\begin{equation}
\begin{aligned}
K(c)&=\operatorname{round}(cN),&
\mathcal S_c&=\{\pi_1,\ldots,\pi_{K(c)}\},\\
R_{\mathrm{unc}}(c)&=\frac{1}{K(c)}\sum_{i\in\mathcal S_c}e_i.
\end{aligned}
\label{eq:selective-risk}
\end{equation}
We compare this ranking with random rejection and an oracle that retains the samples with the
smallest realized errors:
\begin{equation}
\begin{aligned}
\bar R_{\mathrm{rand}}(c)
&=\frac{1}{B}\sum_{b=1}^{B}\frac{1}{K(c)}
\sum_{i\in\mathcal S_c^{(b,\mathrm{rand})}}e_i,\\
R_{\mathrm{oracle}}(c)&=\frac{1}{K(c)}\sum_{i\in\mathcal O_c}e_i.
\end{aligned}
\label{eq:selective-references}
\end{equation}
The random curve answers whether the apparent improvement is merely caused by evaluating fewer
samples, whereas the oracle gives the best possible ranking after the true errors are known. The
distance between the uncertainty curve and these two references therefore separates practical risk
awareness from unavoidable ranking imperfection.
To test whether predicted uncertainty tracks actual error across quantities with different units,
we order each physical metric by its own uncertainty and form percentile bins \(\mathcal B_{m,j}\).
The bin-wise error for metric \(j\) is normalized by that metric's full-coverage MAE, and Panel (b)
reports the mean over the four principal continuous metrics:
\begin{equation}
\widetilde e_{m,j}=
100\frac{\frac{1}{|\mathcal B_{m,j}|}\sum_{i\in\mathcal B_{m,j}}e_{i,j}}
{\operatorname{MAE}^{(j)}_{\mathrm{full}}},
\qquad
\bar e_m^{\mathrm{all}}=\frac{1}{4}\sum_{j=1}^{4}\widetilde e_{m,j}.
\label{eq:uncertainty-bins}
\end{equation}
We use the 0--50, 50--60, \ldots, 90--100 percentile ranges; moving right therefore means that
the model considers the samples progressively less reliable. A value of 100\% equals the average
full-coverage error for the corresponding metric before cross-metric averaging.
Finally, for a rejection fraction \(q=1-c\), the retained-set error and relative improvement are
\begin{equation}
\begin{aligned}
\operatorname{MAE}_{\mathrm{unc}}(q)&=R_{\mathrm{unc}}(1-q),\\
\Delta_{\mathrm{MAE}}(q)&=
\frac{\operatorname{MAE}_{\mathrm{full}}-
\operatorname{MAE}_{\mathrm{unc}}(q)}{\operatorname{MAE}_{\mathrm{full}}}.
\end{aligned}
\label{eq:rejection-improvement}
\end{equation}
Thus, \(R_{\mathrm{unc}}(1)=\operatorname{MAE}_{\mathrm{full}}\), and rejection changes only which
reports are emitted automatically; it does not retrain or alter the underlying predictor.

\begin{figure}[H]
\centering
\begin{minipage}[t]{0.49\columnwidth}
\centering
\includegraphics[width=\linewidth]{figures/selective_prediction_panel_a.png}
\end{minipage}\hfill
\begin{minipage}[t]{0.49\columnwidth}
\centering
\includegraphics[width=\linewidth]{figures/selective_prediction_panel_b.png}
\end{minipage}
\vspace{2pt}
\includegraphics[width=\columnwidth]{figures/selective_prediction_panel_cd.png}
\caption{Reliability-aware selective-prediction diagnostics. (a) Cross-metric normalized
risk--coverage comparison at discrete coverage levels for ensemble-uncertainty ranking,
oracle ranking, and random rejection. (b) Cross-metric normalized MAE versus
uncertainty-percentile bin; 100\% denotes the full-coverage MAE. (c) Per-metric
normalized MAE for first-path delay, angle, power, and $K$-factor under uncertainty
ranking. (d) The same per-metric normalized MAE under oracle ranking. All values are
normalized by the corresponding full-coverage MAE, and lower is better.}
\label{fig:nlos-selective-prediction}
\end{figure}

Figure~\ref{fig:nlos-selective-prediction} shows that ensemble disagreement usefully ranks prediction
risk across heterogeneous physical quantities. After normalizing each metric by its full-coverage
MAE and averaging delay, angle, power, and \(K\)-factor, deferring the highest-uncertainty 10\% and
20\% reduces retained-set error to about 78\% and 69\% of the full-coverage baseline, while random
rejection remains at 100\% in expectation. The bin-wise view is consistent: normalized error rises
from roughly 51\% for the least-uncertain half to 296\% for the highest-uncertainty decile. Although
the remaining gap to the oracle shows that disagreement is not perfectly calibrated, the improvement
over random rejection supports a practical policy: automatically emit low-uncertainty reports and
route deferred samples to additional sensing, paired-numerology fusion, or review. Panel (d) further
shows that the gain is strongest for delay and angle, remains useful for \(K\)-factor, and is weaker
for first-path power.
\subsection{Measured SISO Validation}

We evaluate the pipeline on the Wi-World RealData BUPT mobile-measurement scene
TXH1\_RX01. The acquisition settings in Table~\ref{tab:bupt-config} follow the official dataset
documentation~\cite{wiworldrealdata}. The receiver moves slowly along the measurement route, and
each MAT file covers approximately 0.1 s of a \(1\times1\) SISO link.

\begin{table}[H]
\centering
\caption{Measured-data acquisition, sampling hierarchy, and model input. The independent evaluation
unit is a MAT file; snapshots from the same MAT share one short-time measurement interval.}
\label{tab:bupt-config}
\scriptsize
\setlength{\tabcolsep}{3.0pt}
\renewcommand{\arraystretch}{1.05}
\textbf{Panel A: acquisition configuration}\par\vspace{2pt}
\begin{tabular}{p{0.34\columnwidth}p{0.59\columnwidth}}
\toprule
\textbf{Item} & \textbf{Setting} \\
\midrule
Scene & BUPT TXH1\_RX01; slowly moving receiver \\
Carrier / bandwidth & 3.7 GHz / 200 MHz \\
CIR sampling & 200 MHz; 5 ns delay interval \\
Sounding sequence & PN11; length 2047; 100 MSample/s \\
CIR record & 4094 taps; adjacent CIR interval 20.47 \(\mu\)s \\
MAT contents & Approximately 0.1 s; 4885 or 4886 CIRs \\
Transmitter & LB-20180 horn; 10.206 dBi gain \\
Receiver & SZ-3004000/P omnidirectional bicone; 0.952 dBi gain \\
Channel structure & \(1\times1\) SISO \\
\bottomrule
\end{tabular}

\vspace{4pt}
\textbf{Panel B: data hierarchy}\par\vspace{2pt}
\resizebox{\columnwidth}{!}{%
\begin{tabular}{lcl}
\toprule
\textbf{Level} & \textbf{Amount / size} & \textbf{Role} \\
\midrule
Complete route & \(\sim2997\) MAT / \(\sim900\) GB & Original data pool \\
Local medium pool & 400 MAT / 114.72 GiB & Complete MAT and context retained \\
Debug subset & 10 MAT / \(\sim2.9\) GiB & Pipeline verification; contained in the 60 MAT \\
Formal measured subset & 60 MAT / 17.21 GiB & 60 independent evaluation units \\
Model snapshots & \(60\times32=1920\) & Correlated within each MAT \\
\bottomrule
\end{tabular}}
\end{table}

The 60 formal MAT files are uniformly selected in time from the 400-file local pool. They span
17:45:09.200 to 17:49:58.300 on May 30, 2026, for a total of 289.1 s; the median interval between
selected MAT files is 4.7 s. Each MAT is treated as one approximate measurement location and one
independent evaluation unit. The 10 debug MAT files are included in these 60 files and therefore
must not be counted as 10 additional locations.

Each MAT contains approximately 4885 consecutive CIRs. We uniformly sample 32 snapshots over its
full 0.1 s interval, giving
\begin{equation}
N_{\mathrm{input}}=60\ \mathrm{MAT}\times32\
\frac{\mathrm{snapshots}}{\mathrm{MAT}}=1920\ \mathrm{CSI\ snapshots}.
\label{eq:measured-snapshot-count}
\end{equation}
The 1920 model inputs are not 1920 statistically independent measurement points: the 32 snapshots
within a MAT share the same short-time environment. Each snapshot is converted to 256 frequency
points over 200 MHz, with 781.25 kHz spacing, and represented by real and imaginary channels with
shape \((1,2,256)\). A complete MAT therefore produces an input tensor of shape
\((32,1,2,256)\), with a 1.28 \(\mu\)s unambiguous delay window.

The preprocessing pipeline reads the complex CIR, samples the 32 snapshots, computes their shared
mean power-delay profile (PDP), estimates the noise floor, and constructs one shared soft path mask.
The masked complex CIR is transformed by FFT, resampled to the 256-point frequency grid, corrected
for the fixed system-delay linear phase, and normalized by a per-snapshot common phase before forming
the SISO input. The shared mask threshold is
\begin{equation}
P_{\mathrm{th}}=
\max\!\left(P_{\mathrm{peak}}-25\ \mathrm{dB},
P_{\mathrm{noise}}+6\ \mathrm{dB}\right).
\label{eq:measured-soft-mask}
\end{equation}
Tap weights are zero below the noise floor, one above \(P_{\mathrm{th}}\), and smoothly interpolated
between these limits. Using the same mask for all 32 snapshots in a MAT prevents weak paths from
appearing or disappearing solely because of snapshot-level noise.

The detected first arrival contains both propagation delay and a fixed acquisition-system delay:
\begin{equation}
\tau_{\mathrm{detected}}=
\tau_{\mathrm{system}}+\tau_{\mathrm{propagation}}.
\label{eq:measured-system-delay}
\end{equation}
GNSS positions are successfully matched for 58 of the 60 MAT files. From these, 12 lower-envelope
geometric-LoS candidates are selected, and their median gives
\(\hat{\tau}_{\mathrm{system}}=6199.06\) ns, equivalent to 1239.81 samples at 200 MHz.
Across the candidates, the standard deviation is 37.53 ns, the median absolute deviation is
33.43 ns, and the range is 6161.30--6269.06 ns. We define the reference as the
\emph{PDP-derived calibrated first-arrival delay}
\begingroup
\setlength{\abovedisplayskip}{4pt}
\setlength{\belowdisplayskip}{4pt}
\setlength{\abovedisplayshortskip}{4pt}
\setlength{\belowdisplayshortskip}{4pt}
\begin{equation}
\widetilde{\tau}_{\mathrm{FA}}=
n_{\mathrm{first}}\times5\ \mathrm{ns}-6199.06\ \mathrm{ns}.
\label{eq:measured-calibrated-first-arrival}
\end{equation}
\endgroup
This wording is deliberate: the offset is estimated from automatically selected geometric-LoS
candidates rather than from an independent hardware calibration, so
\(\widetilde{\tau}_{\mathrm{FA}}\) is not strict ground truth.

Only two outputs have measurement-derived references in this experiment:
\(\widetilde{\tau}_{\mathrm{FA}}\) and RMS delay spread \(\tau_{\mathrm{RMS}}\), the latter obtained
from a classical PDP estimator. Other structured outputs, including absolute path power, physical ray
count, interaction count, and receiver-side angle, have no independent reference labels in this
measured subset. They may be generated by the model, but they are excluded from quantitative
accuracy claims. Table~\ref{tab:bupt-measured} therefore evaluates only the two supported delay
quantities.

System-delay calibration reduces first-arrival-delay MAE from 244.78 to 55.25 ns (77.4\%), reduces
the absolute bias from 244.78 to 51.22 ns, and raises Pearson correlation from 0.929 to 0.963. The
remaining negative bias indicates systematic underestimation, although the high correlation shows
that sample-to-sample variation is preserved. Importantly, the residual 55.25 ns is an end-to-end
discrepancy against an automatically constructed reference, rather than pure model error. In signed
form, it contains model, GNSS, first-path-detection, and system-delay-calibration terms. The median
reported RX-position standard deviations are approximately 8.2 m, 11.4 m, and 9.7 m along the
latitude, longitude, and height coordinates; as an order-of-magnitude reference, a 10 m range error
corresponds to about 33 ns of propagation delay. A weak physical first path may also fall below the
PDP detector threshold, causing a later reflection tens of nanoseconds away to be used as the
reference. Finally, the 12 calibration candidates themselves have a 33.43 ns MAD and span
6161.30--6269.06 ns. GNSS uncertainty also propagates into this offset estimate, so these terms are
not independent and the full 55.25 ns cannot be attributed to the predictor alone. For
\(\tau_{\mathrm{RMS}}\), soft-mask AWGN fine-tuning attains 15.941 ns MAE; AWGN training from scratch
gives the best overall result of 13.442 ns MAE, 16.008 ns RMSE, and \(r=0.891\).

\begin{table}[H]
\centering
\caption{Measured SISO delay results on BUPT. Panel A evaluates system-delay calibration for the
PDP-derived first-arrival reference; Panel B compares \(\tau_{\mathrm{RMS}}\) training and
preprocessing. Lower is better for MAE and RMSE; higher is better for Pearson \(r\).}
\label{tab:bupt-measured}
\scriptsize
\setlength{\tabcolsep}{3.0pt}
\renewcommand{\arraystretch}{1.08}
\textbf{Panel A: PDP-derived calibrated first arrival \(\widetilde{\tau}_{\mathrm{FA}}\)}\par\vspace{2pt}
\begin{tabular}{lcc}
\toprule
\textbf{Metric} & \shortstack{\textbf{Soft mask,}\textbf{uncalibrated}} &
\shortstack{\textbf{Soft mask,}\textbf{calibrated}} \\
\midrule
Reference mean (ns) & 758.02 & 758.02 \\
Prediction mean (ns) & 513.25 & 706.81 \\
MAE (ns) $\downarrow$ & 244.78 & \best{\relgain{55.25}} \\
Bias (ns) & $-244.78$ & \relgain{$-51.22$} \\
Pearson $r$ $\uparrow$ & 0.929 & \best{0.963} \\
\bottomrule
\end{tabular}

\vspace{5pt}
\textbf{Panel B: RMS delay spread \(\tau_{\mathrm{RMS}}\)}\par\vspace{2pt}
\resizebox{\columnwidth}{!}{%
\begin{tabular}{llccccc}
\toprule
\textbf{Training} & \textbf{BUPT input} & \textbf{Pred. mean} &
\textbf{MAE $\downarrow$} & \textbf{RMSE $\downarrow$} & \textbf{Bias} &
\textbf{Pearson $r$ $\uparrow$} \\
\midrule
Clean & Raw & 223.994 & 177.906 & 189.028 & $+177.906$ & 0.410 \\
AWGN fine-tuning & Raw & 141.609 & 95.896 & 115.063 & $+95.521$ & 0.423 \\
Clean & Soft mask & 95.016 & 48.929 & 66.988 & $+48.929$ & 0.548 \\
AWGN fine-tuning & Soft mask & 50.583 & 15.941 & 25.976 & \best{$+4.495$} & 0.563 \\
\rowcolor{modelblue}\textbf{AWGN from scratch} & \textbf{Soft mask} & \textbf{57.597} & \best{\relgain{13.442}} & \best{\relgain{16.008}} & \relgain{$+11.509$} & \best{\textbf{0.891}} \\
\bottomrule
\end{tabular}}
\end{table}

The measured-data experiment should be interpreted as an in-sample validation of preprocessing,
calibration, and delay transfer, not as a leakage-free generalization test. In particular, the
6199.06 ns offset is estimated from candidates drawn from the same 60-MAT evaluation subset.
An independent calibration route or direct hardware-delay measurement is needed before claiming
out-of-sample absolute-delay accuracy.

\subsection{Transfer to Downstream Wireless Tasks}

We further evaluate five downstream communication tasks to determine whether the learned CSI
representation is transferable to practical network decisions rather than being useful only for
channel reporting. The downstream heads receive either CSI features, the pretrained embedding, or the
embedding augmented with predicted propagation attributes. Whenever physics features are used, they are
model predictions rather than ground-truth test labels. The five tasks cover localization,
cyclic-prefix (CP) control, link-quality and outage prediction, MCS adaptation, and transmit-beam
selection.

\subsubsection{Two-Dimensional Localization}

For localization, the transmitter and receiver coordinates define the relative-position target
\((\Delta x,\Delta y)=(x_{\mathrm{RX}}-x_{\mathrm{TX}},
y_{\mathrm{RX}}-y_{\mathrm{TX}})\). A regression head maps the CSI representation to this two-dimensional
offset. This task tests whether the representation retains spatial information that is not explicitly
required by the language report. Table~\ref{tab:downstream-localization} shows that augmenting a
trainable CSI encoder with frozen predicted physics gives the lowest mean and median errors and the
largest fractions within 5 and 10 m. The result indicates that the predicted LoS state, delay, and
angular structure provide useful geometric cues beyond a generic language backbone.

\begin{table}[H]
\centering
\caption{Two-dimensional localization on the common test split. Lower is better for mean and median
error; higher is better for distance-threshold accuracy. Underlining marks the best result. The final
row compares the highlighted model with the best of the three language-model baselines (red:
improvement; blue: degradation); threshold accuracies are reported as percentage-point differences.}
\label{tab:downstream-localization}
\scriptsize
\setlength{\tabcolsep}{2.8pt}
\resizebox{\columnwidth}{!}{%
\begin{tabular}{lcccc}
\toprule
Method & Mean (m) $\downarrow$ & Median (m) $\downarrow$ & $\leq5$ m (\%) $\uparrow$ & $\leq10$ m (\%) $\uparrow$ \\
\midrule
\rowcolor{modelblue}\textbf{Scratch + frozen physics} & \best{\textbf{34.15}} & \best{\textbf{10.67}} & \best{\textbf{34.04}} & \best{\textbf{49.02}} \\
Scratch & 47.07 & 20.13 & 13.36 & 35.32 \\
Pretrained embedding & 65.06 & 54.46 & 0.88 & 4.02 \\
Pretrained embedding + physics & 36.91 & 12.04 & 29.70 & 48.20 \\
\midrule
Qwen3.5-2B, full fine-tuning & 100.90 & 101.17 & 0.20 & 0.48 \\
DeepSeek-R1-Qwen3-8B, frozen & 100.79 & 101.09 & 0.16 & 0.50 \\
InternLM3-8B, frozen & 80.07 & 78.95 & 0.86 & 3.10 \\
\midrule
\rowcolor{gainrow}\textbf{Rel. vs. best LLM} & \relgain{57.3\%} & \relgain{86.5\%} & \relgain{33.18 pp} & \relgain{45.92 pp} \\
\bottomrule
\end{tabular}}
\end{table}

\subsubsection{Adaptive Cyclic-Prefix Selection}

For each link, offline ray delays and powers define the shortest candidate CP whose protection
interval contains at least 99\% of the multipath energy. A classifier then selects among the
candidate CP durations from the CSI representation and predicted multipath attributes. The objective
is to avoid an unsafe short CP while also avoiding excessive overhead from an unnecessarily long CP.
Table~\ref{tab:downstream-cp} shows a clear tradeoff: pretrained features alone minimize energy
outside the chosen CP, whereas adding predicted physics yields the highest classification accuracy
and the lowest over-selection rate. Thus, the learned representation supports both reliability-aware
and efficiency-aware CP control.

\begin{table}[H]
\centering
\caption{Adaptive CP selection. Under-selection chooses a CP shorter than the oracle, whereas
over-selection chooses a longer CP. Lower outside-energy values are safer. Underlining marks the
best result for each directed metric. The final row compares the highlighted model with the best of
the three language-model baselines (red: improvement; blue: degradation); predicted CP is
descriptive and is therefore omitted from the relative comparison.}
\label{tab:downstream-cp}
\scriptsize
\setlength{\tabcolsep}{2.0pt}
\resizebox{\columnwidth}{!}{%
\begin{tabular}{lcccccc}
\toprule
Method & Acc. $\uparrow$ & Under $\downarrow$ & Over $\downarrow$ & Pred. CP (ns) & Mean out $\downarrow$ & P95 out $\downarrow$ \\
\midrule
\rowcolor{modelblue}\textbf{Pretrained embedding + physics} & \best{\textbf{70.92\%}} & \textbf{15.92\%} & \best{\textbf{13.16\%}} & \textbf{215.17} & \textbf{4.54\%} & \textbf{9.35\%} \\
Scratch & 51.86\% & 27.32\% & 20.82\% & 189.93 & 6.51\% & 55.51\% \\
Pretrained embedding & 58.62\% & \best{12.22\%} & 29.16\% & 283.01 & \best{2.27\%} & \best{5.21\%} \\
Scratch + frozen physics & 67.76\% & 14.90\% & 17.34\% & 221.40 & 3.57\% & 5.39\% \\
\midrule
Qwen3.5-2B, full fine-tuning & 51.76\% & 19.72\% & 28.52\% & 290.20 & 2.80\% & 10.18\% \\
InternLM3-8B, frozen + CSI & 55.58\% & 21.24\% & 23.18\% & 208.43 & 6.24\% & 57.31\% \\
DeepSeek-8B, frozen + CSI & 58.30\% & 15.74\% & 25.96\% & 247.43 & 3.27\% & 11.37\% \\
\midrule
\rowcolor{gainrow}\textbf{Rel. vs. best LLM} & \relgain{21.6\%} & \relloss{1.1\%} & \relgain{43.2\%} & --- & \relloss{62.1\%} & \relgain{8.2\%} \\
\bottomrule
\end{tabular}}
\end{table}

\subsubsection{Link Quality and Outage Prediction}

The link-quality task estimates received power and decides whether a link remains operational under
the specified noise condition. It therefore tests whether the learned power, LoS, and multipath cues
are sufficiently calibrated for reliability-oriented decisions. As shown in
Table~\ref{tab:downstream-outage}, the CSI2TEXT-PGMT total-power head gives the lowest power MAE,
while all methods exceed 98.9\% outage accuracy and F1. The language baselines obtain slightly higher
binary scores but incur larger power errors, showing that an accurate outage decision does not by
itself guarantee a faithful continuous link-quality estimate.

\begin{table}[H]
\centering
\caption{Link-quality and outage prediction. Lower is better for power MAE; higher is better for
outage accuracy and F1. Underlining marks the best result. The final row compares the highlighted
model with the best of the three language-model baselines (red: improvement; blue: degradation).}
\label{tab:downstream-outage}
\scriptsize
\setlength{\tabcolsep}{2.6pt}
\resizebox{\columnwidth}{!}{%
\begin{tabular}{llccc}
\toprule
Method & Training & Power MAE $\downarrow$ & Outage acc. $\uparrow$ & Outage F1 $\uparrow$ \\
\midrule
\rowcolor{modelblue}\textbf{\model{}} & \textbf{Total-power head} & \best{\textbf{2.1098 dB}} & \textbf{98.92\%} & \textbf{98.90\%} \\
DeepSeek-R1-Qwen3-8B & Frozen + adapter/head & 3.0165 dB & 99.52\% & 99.51\% \\
InternLM3-8B-Instruct & Frozen + adapter/head & 2.8207 dB & 99.52\% & 99.51\% \\
Qwen3.5-2B & Full fine-tuning & 2.4603 dB & \best{99.56\%} & \best{99.55\%} \\
\midrule
\rowcolor{gainrow}\multicolumn{2}{l}{\textbf{Rel. vs. best LLM}} & \relgain{14.2\%} & \relloss{0.6\%} & \relloss{0.7\%} \\
\bottomrule
\end{tabular}}
\end{table}

\subsubsection{MCS Link Adaptation}

For MCS adaptation, each sample combines the CSI representation, transmit power, and predicted
propagation features. Effective-SNR and calibrated BLER constraints define the highest feasible MCS,
including an outage class when no MCS satisfies the target. The classifier must therefore maximize
spectral efficiency without violating the reliability constraint. Table~\ref{tab:downstream-mcs}
shows that the scratch CSI2TEXT-PGMT configuration reaches 99.05\% accuracy, only 0.53\% BLER
violations, and 99.90\% of oracle goodput. This demonstrates that the structured CSI interface can
support near-oracle link adaptation without relying on a large language decoder at inference.

\begin{table}[H]
\centering
\caption{MCS link adaptation under calibrated BLER constraints. Lower is better for MCS MAE and
BLER violation; higher is better for accuracy, goodput ratio, and outage recall. The final row
compares the highlighted model with the best of the three language-model baselines (red:
improvement; blue: degradation).}
\label{tab:downstream-mcs}
\scriptsize
\setlength{\tabcolsep}{2.0pt}
\resizebox{\columnwidth}{!}{%
\begin{tabular}{lccccc}
\toprule
Method & Acc. $\uparrow$ & MCS MAE $\downarrow$ & BLER viol. $\downarrow$ & Goodput $\uparrow$ & Outage recall $\uparrow$ \\
\midrule
\rowcolor{modelblue}\textbf{Scratch} & \best{\textbf{99.05\%}} & \best{\textbf{0.0098}} & \best{\textbf{0.53\%}} & \best{\textbf{99.90\%}} & \textbf{99.49\%} \\
Pretrained embedding & 91.80\% & 0.1016 & 3.56\% & 96.67\% & 95.94\% \\
Pretrained embedding + physics & 92.20\% & 0.0948 & 4.40\% & 97.11\% & 95.21\% \\
Scratch + frozen physics & 98.07\% & 0.0321 & 1.38\% & 99.63\% & 98.02\% \\
\midrule
Qwen3.5-2B & 95.44\% & 0.0458 & 1.22\% & 98.54\% & \best{99.59\%} \\
DeepSeek-R1-Qwen3-8B & 94.56\% & 0.0544 & 4.76\% & 98.00\% & 98.29\% \\
InternLM3-8B & 94.98\% & 0.0504 & 2.46\% & 98.60\% & 98.62\% \\
\midrule
\rowcolor{gainrow}\textbf{Rel. vs. best LLM} & \relgain{3.8\%} & \relgain{78.6\%} & \relgain{56.6\%} & \relgain{1.3\%} & \relloss{0.1\%} \\
\bottomrule
\end{tabular}}
\end{table}

\subsubsection{Adaptive Transmit-Beam Search}

The beam task reconstructs the 64 candidate beams of an $8\times8$ UPA and uses the beam with the
largest average received gain as the oracle label. A 64-class head predicts the best transmit-beam
index from the full CSI representation. This evaluates whether the representation preserves the
directional information required for beam management. Table~\ref{tab:downstream-beam} shows that
scratch training augmented with frozen predicted physics gives the strongest overall result: 87.84\%
Top-1 accuracy, 96.52\% Top-3 accuracy, and only 1.33 dB mean gain loss. The improvement is largest
for NLoS links, where indirect multipath cues are particularly valuable.

\begin{table}[H]
\centering
\caption{Full-CSI transmit-beam prediction over 64 UPA beams. Lower is better for gain loss; higher
is better for all other metrics. Underlining marks the best result. The final row compares the
highlighted model with the best of the three language-model baselines (red: improvement; blue:
degradation).}
\label{tab:downstream-beam}
\scriptsize
\setlength{\tabcolsep}{1.8pt}
\resizebox{\columnwidth}{!}{%
\begin{tabular}{lccccccc}
\toprule
Method & Top-1 & Top-3 & $\leq1$ dB & Mean $\downarrow$ & P90 $\downarrow$ & LoS & NLoS \\
\midrule
\rowcolor{modelblue}\textbf{Scratch + frozen physics} & \best{\textbf{87.84\%}} & \best{\textbf{96.52\%}} & \best{\textbf{89.92\%}} & \best{\textbf{1.33}} & \best{\textbf{1.11}} & \textbf{98.80\%} & \best{\textbf{76.88\%}} \\
Pretrained embedding & 87.10\% & 95.40\% & 89.40\% & 1.57 & 1.60 & 98.68\% & 75.52\% \\
Pretrained embedding + physics & 87.02\% & 95.20\% & 89.10\% & 1.62 & 2.05 & 98.80\% & 75.24\% \\
Scratch & 86.04\% & 96.28\% & 88.38\% & 1.65 & 2.75 & \best{98.88\%} & 73.20\% \\
\midrule
Qwen3.5-2B & 77.30\% & 91.10\% & 80.35\% & 2.53 & 10.79 & 92.50\% & 62.10\% \\
InternLM3-8B-Instruct & 77.60\% & 90.15\% & 81.25\% & 2.51 & 10.05 & 92.60\% & 62.60\% \\
DeepSeek-R1-Qwen3-8B & 49.00\% & 76.15\% & 51.65\% & 8.04 & 25.72 & 66.50\% & 31.50\% \\
\midrule
\rowcolor{gainrow}\textbf{Rel. vs. best LLM} & \relgain{13.2\%} & \relgain{5.9\%} & \relgain{10.7\%} & \relgain{47.0\%} & \relgain{89.0\%} & \relgain{6.7\%} & \relgain{22.8\%} \\
\bottomrule
\end{tabular}}
\end{table}

Across the five tasks, no single feature choice dominates every secondary metric, but the results
consistently show that CSI-native representations and predicted propagation attributes provide a
compact and effective interface for network control. In particular, physics assistance is most useful
for geometry- and multipath-sensitive decisions, while lightweight CSI heads can match or exceed
direct language-model baselines at substantially lower inference cost.
